\subsection{W 在诸室的并中的基本域}

我们沿用第 4 号的记号。对于 \(s \in S\) ，用 \(H_s\) 表示 \(E^*\) 中正交于 \(e_s\) 的超平面，用 \(\overline{A_s}\) 表示由 \(x^* \in E^*\) （满足 \(\langle x^*, e_s \rangle \geqslant 0\) ）构成的集合，并用 \(\overline{C}\) 表示诸 \(\overline{A_s}\) （\(s \in S\) ）的交。对于弱拓扑 \(\sigma(E^*, E)\) （它由 \(E^*\) 与 \(E\) 之间的对偶性定义；《拓扑向量空间》，第 II 章，§6，第 2 号），诸 \(\overline{A_s}\) 是闭半空间，而 \(\overline{C}\) 是一个闭凸锥。此外，\(\overline{C}\) 是 C 的闭包；事实上，若 \(x^* \in \overline{C}\) 且 \(y^* \in C\) ，则 \(x^* + ty^* \in C\) （对每个实数 \(t > 0\) ），且 \(x^* = \lim_{t \to 0} (x^* + ty^*)\) 。

对于 \(X\subset S\) ，置

\[
C _ {X} = \left(\bigcap_ {s \in X} H _ {s}\right) \cap \left(\bigcap_ {s \in S - X} A _ {s}\right).
\]

我们有 \(C_X \subset \overline{C}\) ，\(C_{\varnothing} = C\) 且 \(C_S = \{0\}\) 。诸集合 \(C_X\) （\(X \in B(S)\) ）构成 \(\overline{C}\) 的一个划分。

另一方面，回忆（第 IV 章，§1，第 8 号）\(\mathbf{W}_{\mathbf{X}}\) 表示 \(\mathbf{W}\) 的由 \(\mathbf{X}\) 生成的子群。显然 \(w(x^{*}) = x^{*}\) （对于 \(w \in \mathbf{W}_{\mathbf{X}}\) 与 \(x^{*} \in \mathbf{C}_{\mathbf{X}}\) ）。

命题 5. 设 \(X, X' \subset S\) 且 \(w, w' \in W\) 。若 \(w(C_X) \cap w'(C_{X'}) \neq \emptyset\) ，则 \(X = X'\) ，\(wW_X = w'W_{X'}\) 且 \(w(C_X) = w'(C_{X'})\) 。

我们立即化归到 \(w' = 1\) 的情形。证明对长度 \(n\) （\(w\) 的长度）作归纳。若 \(n = 0\) ，断言是显然的。若 \(l(w) > 0\) ，则存在 \(s \in S\) 使 \(l(sw) = l(w) - 1\) ，于是（参见~第 4 号末尾）\(w(C) \subset s(A_s)\) ，从而 \(w(\overline{C}) \subset s(\overline{A_s})\) 。由于 \(\overline{C} \subset \overline{A_s}\) ，由此得到

\[
\overline {{{\mathrm{C}}}} \cap w (\overline {{{\mathrm{C}}}}) \subset \mathrm{H} _ {s}.
\]

于是 \(s(x^{*}) = x^{*}\) （对所有 \(x^{*} \in \overline{\mathrm{C}} \cap w(\overline{\mathrm{C}})\) ），从而更对所有

\[
x ^ {*} \in \mathrm{C} _ {\mathrm{X} ^ {\prime}} \cap w (\mathrm{C} _ {\mathrm{X}}).
\]

因此，关系 \(C_{X'} \cap w(C_X) \neq \varnothing\) 一方面蕴含 \(C_{X'} \cap H_s \neq \varnothing\) ，从而 \(s \in X'\) ；另一方面蕴含 \(C_{X'} \cap sw(C_X) \neq \varnothing\) 。于是归纳假设表明 \(X = X'\) 且 \(swW_X = W_{X'} = W_X\) ，故 \(sw \in W_X\) ，且 \(w \in W_X\) （由于 \(s \in W_X\) ）。由此得到 \(wW_X = W_{X'}\) 以及 \(w(C_X) = C_X = C_{X'}\) 。

推论. 设 X 是 S 的子集，\(x^{*}\) 是 \(C_{X}\) 的一个元素。\(x^{*}\) 在 W 中的稳定子是 \(W_{X}\) 。

现在设 U 为诸 \(w(\overline{\mathbf{C}})\) （\(w \in W\) ）的并，并设 F 为 U 中形如 \(w(\mathrm{C}_{X})\) 的子集（其中 \(X \subset S\) 且 \(w \in W\) ）所成的集合。由上所述，F 是 U 的一个划分。

命题 6. (i) 锥 U 是凸的。

\begin{enumerate}
\def\labelenumi{(\roman{enumi})}
\setcounter{enumi}{1}
\item
  U 中所含的每个闭线段只与 \(\mathfrak{F}\) 中的有限多个元素相交。
\item
  锥 \(\overline{\mathbf{C}}\) 是 W 在 U 上作用的一个基本域。
\end{enumerate}

为证明 (iii)，只需证明：若 \(x^{*}, y^{*} \in \overline{\mathbb{C}}\) 与 \(w \in \mathbb{W}\) 满足 \(w(x^{*}) = y^{*}\) ，则 \(x^{*} = y^{*}\) 。现在存在两个子集 \(\mathbf{X}\) 与 \(\mathbf{Y}\) （\(\mathbf{S}\) 中的），使得 \(x^{*} \in \mathrm{C}_{\mathbf{X}}\) 且 \(y^{*} \in \mathrm{C}_{\mathbf{Y}}\) ；我们有 \(w(\mathrm{C}_{\mathbf{X}}) \cap \mathrm{C}_{\mathbf{Y}} \neq \emptyset\) ，而命题 5 表明 \(\mathbf{X} = \mathbf{Y}\) 且 \(w \in \mathrm{W}_{\mathbf{X}}\) ，这就蕴含 \(x^{*} = y^{*}\) 。

现在设 \(x^{*}, y^{*} \in \mathrm{U}\) ；我们将证明线段 \([x^{*}y^{*}]\) 被 \(\mathfrak{F}\) 中有限多个元素所覆盖，这就同时建立了 (i) 与 (ii)。把 \(x^{*}\) 与 \(y^{*}\) 用 \(\mathrm{W}\) 的同一个元素来变换，我们可以假设 \(x^{*} \in \overline{\mathrm{C}}\) 。设 \(w \in \mathrm{W}\) 满足 \(y^{*} \in w(\overline{\mathrm{C}})\) 。我们对 \(w\) 的长度作归纳论证。对于 \(s \in \mathrm{S}\) ，关系 \(w(\overline{\mathrm{C}}) \not\subset \overline{\mathrm{A}_{s}}\) 等价于 \(w(\mathrm{C}) \notin \mathrm{A}_{s}\) ，从而等价于 \(l(sw) < l(w)\) （参见~第 4 号）。于是第 IV 章 §1 第 8 号的命题 7 表明，只存在有限多个 \(s \in \mathrm{S}\) 满足 \(w(\overline{\mathrm{C}}) \not\subset \overline{\mathrm{A}_{s}}\) 。于是，集合 \(T\) （由那些 \(s \in S\) 组成，其中 \(\langle y^{*}, e_{s} \rangle < 0\) ）是有限的。另一方面，交 \(\overline{\mathrm{C}} \cap [x^{*}y^{*}]\) 是一个闭线段 \([x^{*}z^{*}]\) 。若 \(z^{*} = y^{*}\) ，即若 \(y^{*} \in \overline{\mathrm{C}}\) ，则存在子集 \(X\) 与 \(Y\) （\(S\) 中的），使得 \(x^{*} \in C_{X}\) 且 \(y^{*} \in C_{Y}\) 。开线段 \(|x^{*}y^{*}|\) 于是含于 \(C_{X \cap Y}\) ，故 \([x^{*}y^{*}] \subset C_{X} \cup C_{Y} \cup C_{X \cap Y}\) 。若 \(z^{*} \neq y^{*}\) ，则存在 \(s \in T\) 使 \(z^{*} \in H_{s}\) 。于是 \(w(C) \not\subset A_{s}\) 且 \(l(sw) < l(w)\) 。从而归纳假设表明，线段 \([z^{*}y^{*}] = s([z^{*}(s(y^{*}))])\) 被 \(\mathfrak{F}\) 的有限多个元素所覆盖，因而

\[
\left[ x ^ {*} y ^ {*} \right] = \left[ x ^ {*} z ^ {*} \right] \cup \left[ z ^ {*} y ^ {*} \right]
\]

因为 \([x^{*}z^{*}]\subset \overline{\mathbf{C}}.\)

